% NOVA overview + RelTwin detail. Standalone design preview, 2026-09-15.
% Top timelines and excerpt layout are schematic; panel (c) uses unchanged
% measured intervals and IoUs from figures/overview_data.tex.
% Solid arrows: inference/data flow. Purple dashed arrow: training influence.
% Candidate families and configuration-dependent routing follow sec:system.
% The displayed RelTwin loss is the three-term controlled training objective;
% the system's consistency/SetPO continuation is specified in the paper text.
% Generated from frozen cached predictions; see notes/figure_evidence.json.
\def\FigGTPlusStart{1.5}
\def\FigGTPlusEnd{11.75}
\def\FigGTMinusStart{14.75}
\def\FigGTMinusEnd{25.0}
\def\FigSFTPlusStart{14.75}
\def\FigSFTPlusEnd{24.75}
\def\FigSFTMinusStart{14.75}
\def\FigSFTMinusEnd{24.75}
\def\FigCandPlusStart{1.6}
\def\FigCandPlusEnd{8.7}
\def\FigCandMinusStart{14.75}
\def\FigCandMinusEnd{24.75}
\def\FigSftPlusIoU{0.00}
\def\FigSftMinusIoU{0.98}
\def\FigCandPlusIoU{0.69}
\def\FigCandMinusIoU{0.98}
\def\FigSameBefore{37}
\def\FigJointBefore{53.13}
\def\FigSameAfter{1}
\def\FigJointAfter{84.38}

\definecolor{nvInk}{HTML}{242424}
\definecolor{nvCanvas}{HTML}{FBF8F3}
\definecolor{nvCyan}{HTML}{658C99}
\definecolor{nvCyanFill}{HTML}{BFDCE7}
\definecolor{nvCyanBg}{HTML}{DFEDF2}
\definecolor{nvPurple}{HTML}{726082}
\definecolor{nvPurpleFill}{HTML}{CAC2D7}
\definecolor{nvPurpleBg}{HTML}{DDD6E5}
\definecolor{nvGreen}{HTML}{5D7947}
\definecolor{nvGreenFill}{HTML}{D9E4C2}
\definecolor{nvGreenBar}{HTML}{A8BF8C}
\definecolor{nvPurpleBar}{HTML}{B5A3C9}
\definecolor{nvApricot}{HTML}{F5D8BF}
\definecolor{nvApricotLine}{HTML}{A67F5D}
\begin{tikzpicture}[x=1cm,y=1cm,
 every node/.style={font=\fontsize{9}{10.3}\selectfont,text=nvInk,inner sep=1.4pt},
 flow/.style={-{Latex[length=1.5mm]},line width=.7pt,draw=nvInk},
 train/.style={-{Latex[length=1.6mm]},dash pattern=on 3pt off 2pt,line width=.9pt,draw=nvPurple},
 heading/.style={font=\fontsize{10}{11.3}\selectfont\bfseries},
 panel/.style={font=\fontsize{10}{11.3}\selectfont\bfseries},
 box/.style={draw=nvCyan,line width=.7pt,fill=white}]
\path[use as bounding box] (0,0) rectangle (17.4,9.70);
\path[fill=nvCanvas] (0,0) rectangle (17.4,9.70);

% System identity and inference lane.
\node[anchor=west,font=\fontsize{14}{15}\selectfont\bfseries] at (.16,9.38) {NOVA};
\node[anchor=west,font=\fontsize{10}{11.3}\selectfont] at (1.97,9.38) {Query-specific audio grounding};
\node[anchor=east] at (17.20,9.38) {System overview};

% Input block.
\draw[box,fill=nvCyanBg!65] (.16,7.17) rectangle (2.45,8.99);
\node[heading] at (1.305,8.71) {Audio + query};
\draw[nvCyan!40,line width=.5pt] (.44,8.20)--(2.17,8.20);
\foreach \xx/\hh in {.50/.06,.62/.14,.74/.23,.86/.10,.98/.18,1.10/.29,1.22/.16,1.34/.07,1.46/.13,1.58/.23,1.70/.15,1.82/.06,1.94/.12,2.06/.04}{
 \draw[nvCyan,line width=.9pt] (\xx,{8.20-\hh})--(\xx,{8.20+\hh});
}
\node[align=center] at (1.305,7.55) {``dog followed\\by rooster''};

% Candidate generation; model identities are checkpoint families, not
% additional inference-time training operators.
\draw[box,fill=nvCyanBg!35] (3.04,7.17) rectangle (6.18,8.99);
\path[fill=nvCyanFill] (3.045,8.43) rectangle (6.175,8.985);
\node[font=\fontsize{9}{10.3}\selectfont\bfseries] at (4.61,8.71) {Candidate generation};
\node at (4.61,8.15) {Official / SFT / SetPO};
\foreach \yy/\aa/\bb in {7.87/3.50/4.36,7.66/4.49/5.77,7.45/3.58/4.52}{
 \draw[nvCyan!35,line width=.45pt] (3.46,\yy)--(5.78,\yy);
 \draw[draw=nvCyan,fill=nvCyanFill,line width=.5pt] (\aa,{\yy-.06}) rectangle (\bb,{\yy+.06});
}

% Evidence routing: keep and drop illustrate the same candidate mask.
\draw[box,draw=nvGreen,fill=nvGreenFill!25] (6.78,7.17) rectangle (10.07,8.99);
\node[heading] at (8.425,8.71) {Evidence routing*};
\node[anchor=west] at (6.96,8.20) {Keep};
\node[anchor=west] at (6.96,7.80) {Drop};
\foreach \yy in {8.20,7.80}{
 \draw[nvGreen!35,line width=.5pt] (7.92,\yy)--(9.80,\yy);
}
\draw[draw=nvGreen,fill=nvGreenBar,line width=.5pt] (8.48,8.115) rectangle (9.19,8.285);
\draw[draw=nvGreen,fill=nvGreenFill,line width=.5pt] (7.92,7.715) rectangle (8.48,7.885);
\draw[draw=nvGreen,fill=nvGreenFill,line width=.5pt] (9.19,7.715) rectangle (9.80,7.885);
\draw[nvGreen,dash pattern=on 1pt off 1pt,line width=.5pt] (8.48,7.715) rectangle (9.19,7.885);
\node at (8.425,7.43) {Query-existence scores};

% Boundary alignment, with the original outline and a refined filled span.
\draw[box,draw=nvApricotLine,fill=nvApricot!22] (10.67,7.17) rectangle (13.99,8.99);
\node[heading] at (12.33,8.71) {Boundary refinement};
\node at (12.33,8.22) {Duration-scaled radius};
\draw[nvApricotLine!40,line width=.5pt] (11.09,7.84)--(13.58,7.84);
\draw[draw=nvApricotLine,fill=nvApricot,line width=.7pt] (11.52,7.76) rectangle (13.19,7.92);
\draw[nvApricotLine,dash pattern=on 2pt off 1pt,line width=.7pt] (11.73,7.69) rectangle (13.03,7.99);
\draw[{Latex[length=.9mm]}-{Latex[length=.9mm]},nvApricotLine,line width=.6pt] (11.38,8.06)--(11.88,8.06);
\draw[{Latex[length=.9mm]}-{Latex[length=.9mm]},nvApricotLine,line width=.6pt] (12.86,8.06)--(13.35,8.06);
\node at (12.33,7.43) {DP + bootstrap gate};

% Output intervals (schematic, no claimed example prediction).
\draw[box,fill=nvCyanBg!40] (14.59,7.17) rectangle (17.24,8.99);
\node[heading] at (15.915,8.71) {Final intervals};
\draw[flow,line width=.5pt] (14.89,8.18)--(16.98,8.18);
\draw[draw=nvPurple,fill=nvPurpleBar,line width=.65pt] (15.22,8.08) rectangle (16.39,8.28);
\node at (15.915,7.72) {$\widehat S=\{[\hat s_k,\hat e_k]\}_k$};
\node at (15.915,7.40) {Localized events};

\foreach \xa/\xb in {2.45/3.04,6.18/6.78,10.07/10.67,13.99/14.59}{
 \draw[flow] (\xa,8.10)--(\xb,8.10);
}
\node[anchor=east] at (17.20,6.79) {*Evidence routing is configuration-dependent.};

% Focus area: visibly connected to candidate training, not inserted into the
% runtime data-flow lane. No directed connection to the output example.
\draw[draw=nvPurple!65,fill=nvPurpleBg!12,line width=.75pt] (.16,.17) rectangle (17.24,6.16);
\path[fill=nvPurpleFill!65] (.165,5.65) rectangle (17.235,6.155);
\node[anchor=west,font=\fontsize{11}{12.5}\selectfont\bfseries] at (.37,5.91) {RelTwin};
\node[anchor=west,font=\fontsize{10}{11.3}\selectfont] at (2.02,5.91) {Locally valid timestamp contrasts};
\node[anchor=east] at (17.03,5.91) {Training + diagnostic example};
\draw[train] (4.61,6.16)--(4.61,7.16);
\node[anchor=west,fill=nvCanvas] at (4.80,6.52) {Trains a candidate generator};

% Subpanel dividers and headings.
\draw[nvPurple!24,line width=.55pt] (5.37,.48)--(5.37,5.44);
\draw[nvPurple!24,line width=.55pt] (11.20,.48)--(11.20,5.44);
\node[panel] at (2.72,5.36) {(a) Inverse-relation pairs};
\node[panel] at (8.285,5.36) {(b) Answer competition};
\node[panel] at (14.22,5.36) {(c) Query-specific localization};

% (a): explicit relation pairing within one recording.
\node at (2.72,4.96) {One audio $x$, shared $A$ and $B$};
\draw[nvCyan!35,line width=.5pt] (.49,4.43)--(4.99,4.43);
\foreach \xa/\xb/\event/\eventfill/\eventstroke in {
 .56/1.34/A/nvCyanFill/nvCyan,1.46/2.24/B/nvApricot/nvApricotLine,
 3.03/3.81/B/nvApricot/nvApricotLine,3.93/4.71/A/nvCyanFill/nvCyan}{
 \draw[draw=\eventstroke,fill=\eventfill,line width=.65pt] (\xa,4.26) rectangle (\xb,4.62);
 \node at ({(\xa+\xb)/2},4.44) {$\event$};
}
\draw[flow,line width=.55pt] (.49,4.09)--(4.99,4.09);
\node[anchor=west] at (4.96,4.10) {$t$};
\draw[nvGreen,line width=.85pt] (.56,3.99)--(.56,3.86)--(2.24,3.86)--(2.24,3.99);
\draw[nvPurple,line width=.85pt] (3.03,3.99)--(3.03,3.86)--(4.71,3.86)--(4.71,3.99);
\node at (1.40,3.66) {$W_+: A\rightarrow B$};
\node at (3.87,3.66) {$W_-: B\rightarrow A$};
\path[fill=nvGreenFill] (.48,2.86) rectangle (5.04,3.35);
\draw[nvGreen,line width=1.2pt] (.48,2.86)--(.48,3.35);
\node[anchor=west] at (.62,3.105) {$q_+$: dog followed by rooster};
\path[fill=nvPurpleFill] (.48,2.17) rectangle (5.04,2.66);
\draw[nvPurple,line width=1.2pt] (.48,2.17)--(.48,2.66);
\node[anchor=west] at (.62,2.415) {$q_-$: rooster followed by dog};
\node[align=center] at (2.72,1.54) {Both windows have acoustic support.\\The query selects its window.};
\node at (2.72,.78) {Inverse query $\Longrightarrow$ exchanged labels};

% (b): frozen backbone, four query/answer scores, matched diagonal labels.
\draw[box,fill=nvCyanBg] (5.72,4.40) rectangle (10.88,5.04);
\node[align=center] at (8.30,4.72) {Shared audio-language model\\Frozen backbone + trainable LoRA};
% The input bus stays clear of the panel divider and the target brackets.
\draw[nvInk,line width=.65pt] (5.04,3.105)--(5.22,3.105);
\draw[nvInk,line width=.65pt] (5.04,2.415)--(5.22,2.415)--(5.22,4.72);
\fill[nvInk] (5.22,3.105) circle (.8pt);
\draw[flow] (5.22,4.72)--(5.72,4.72);
\node[anchor=south] at (5.45,4.78) {$x,q_{\pm}$};
\draw[flow] (8.30,4.40)--(8.30,4.12);
\node at (7.72,3.98) {$y(W_+)$};
\node at (9.38,3.98) {$y(W_-)$};
\node at (6.57,3.55) {$q_+$};
\node at (6.57,2.95) {$q_-$};
\path[fill=nvApricot] (6.89,3.25) rectangle (8.55,3.85);
\path[fill=nvApricot] (8.55,2.65) rectangle (10.21,3.25);
\draw[nvPurple,line width=.7pt] (6.89,2.65) rectangle (10.21,3.85);
\draw[nvPurple!70,line width=.55pt] (8.55,2.65)--(8.55,3.85);
\draw[nvPurple!70,line width=.55pt] (6.89,3.25)--(10.21,3.25);
\node at (7.72,3.55) {$s_{++}\ \uparrow$};
\node at (9.38,3.55) {$s_{+-}\ \downarrow$};
\node at (7.72,2.95) {$s_{-+}\ \downarrow$};
\node at (9.38,2.95) {$s_{--}\ \uparrow$};
\node at (8.30,2.36) {$s$: mean token log-likelihood};
\draw[flow] (8.30,2.17)--(8.30,1.94);
\node at (8.30,1.69) {$\mathcal L_{\rm cand}=-\frac12[\log p_+(+)+\log p_-(-)]$};
\node at (8.30,1.21) {Query-conditioned answer preference};
\node at (8.30,.74) {$\mathcal L=\mathcal L_{\rm seq}+\lambda\mathcal L_{\rm cand}+\beta\mathcal L_{\rm replay}$};

% (c): identical data to the current paper. This is a RelTwin-vs-SFT diagnostic,
% not an example of the full routed/refined NOVA system.
\foreach \xl/\xr in {12.09/14.43,14.76/17.10}{
 \path[fill=nvCyanBg] (\xl,4.63) rectangle (\xr,5.04);
}
\node[heading] at (13.26,4.835) {SFT};
\node[heading] at (15.93,4.835) {RelTwin};
\foreach \xbase in {12.18,14.85}{
 \draw[nvInk,line width=.55pt] (\xbase,4.29)--({\xbase+2.14},4.29);
 \foreach \sec/\lab in {0/0,10/10,20/20,27/27}{
  \pgfmathsetmacro{\xx}{\xbase+2.14*\sec/27}
  \draw[nvInk,line width=.5pt] (\xx,4.25)--(\xx,4.33);
  \node[anchor=south] at (\xx,4.33) {\lab};
 }
}
\node[anchor=east] at (11.96,3.38) {$q_+$};
\node[anchor=east] at (11.96,2.29) {$q_-$};
\foreach \xbase/\ystart/\ystop/\pstart/\pstop/\yy/\barstroke/\barfill in {
 12.18/\FigGTPlusStart/\FigGTPlusEnd/\FigSFTPlusStart/\FigSFTPlusEnd/3.38/nvGreen/nvGreenBar,
 14.85/\FigGTPlusStart/\FigGTPlusEnd/\FigCandPlusStart/\FigCandPlusEnd/3.38/nvGreen/nvGreenBar,
 12.18/\FigGTMinusStart/\FigGTMinusEnd/\FigSFTMinusStart/\FigSFTMinusEnd/2.29/nvPurple/nvPurpleBar,
 14.85/\FigGTMinusStart/\FigGTMinusEnd/\FigCandMinusStart/\FigCandMinusEnd/2.29/nvPurple/nvPurpleBar}{
 \draw[\barstroke!28,line width=.45pt] (\xbase,\yy)--({\xbase+2.14},\yy);
 \pgfmathsetmacro{\predstart}{\xbase+2.14*\pstart/27}
 \pgfmathsetmacro{\predstop}{\xbase+2.14*\pstop/27}
 \draw[draw=\barstroke,fill=\barfill,line width=.55pt]
  (\predstart,{\yy-.09}) rectangle (\predstop,{\yy+.09});
 \pgfmathsetmacro{\targetstart}{\xbase+2.14*\ystart/27}
 \pgfmathsetmacro{\targetstop}{\xbase+2.14*\ystop/27}
 \draw[draw=\barstroke,fill=none,line width=.9pt]
  (\targetstart,{\yy-.15}) rectangle (\targetstop,{\yy+.15});
}
\node[anchor=east] at (14.30,3.73) {IoU \FigSftPlusIoU};
\node[anchor=east,font=\fontsize{9}{10.3}\selectfont\bfseries] at (16.97,3.73) {IoU \FigCandPlusIoU};
\node[anchor=east] at (14.30,2.64) {IoU \FigSftMinusIoU};
\node[anchor=east] at (16.97,2.64) {IoU \FigCandMinusIoU};
\draw[draw=nvInk,fill=none,line width=.8pt] (12.20,1.66) rectangle (12.55,1.84);
\node[anchor=west] at (12.65,1.75) {target};
\draw[draw=nvInk,fill=nvPurpleBar,line width=.55pt] (14.33,1.66) rectangle (14.68,1.84);
\node[anchor=west] at (14.78,1.75) {prediction};
\node[align=center] at (14.25,.98) {Relation-development example:\\shared time axis in seconds.};
\end{tikzpicture}
